\DocumentMetadata{
  pdfversion=2.0,
  pdfstandard=ua-2,
  lang=en-US,
  tagging=on,
  tagging-setup={math/setup=mathml-SE,tabsorder=structure}
}
\documentclass[11pt,letterpaper]{article}
\usepackage[margin=0.68in,includefoot]{geometry}
\usepackage{newcomputermodern}
\usepackage{amsmath,amscd,mathtools}
\providecommand{\square}{\mdlgwhtsquare}
\usepackage[hidelinks]{hyperref}
\hypersetup{
  pdftitle={Variational Analysis/Numerical Optimization PhD Exam, May 2025},
  pdfauthor={University of Florida Department of Mathematics},
  pdfsubject={Graduate mathematics examination},
  pdfdisplaydoctitle=true,
  bookmarksopen=true
}
\setcounter{secnumdepth}{0}
\setlength{\parindent}{0pt}
\setlength{\parskip}{0.28em}
\setlength{\emergencystretch}{3em}
\setlength{\leftmargini}{2.15em}
\setlength{\leftmarginii}{2.65em}
\setlength{\leftmarginiii}{3em}
\renewcommand{\labelenumi}{\textbf{\theenumi.}}
\pagestyle{empty}
\raggedbottom
\allowdisplaybreaks
\newenvironment{examproblems}[1]{%
  \begin{enumerate}%
  \setcounter{enumi}{#1}%
  \setlength{\topsep}{0.35em}%
  \setlength{\partopsep}{0pt}%
  \setlength{\itemsep}{0.48em}%
  \setlength{\parsep}{0.18em}%
}{\end{enumerate}}
\newenvironment{examsubparts}{%
  \begin{enumerate}%
  \setlength{\topsep}{0.18em}%
  \setlength{\partopsep}{0pt}%
  \setlength{\itemsep}{0.16em}%
  \setlength{\parsep}{0.1em}%
}{\end{enumerate}}
\newenvironment{examinstructions}{%
  \par\begingroup\itshape
}{%
  \par\endgroup\vspace{0.18em}
}
\makeatletter
\renewcommand\section{\@startsection{section}{1}{0pt}%
  {-1.25ex plus -.25ex minus -.1ex}%
  {0.55ex plus .1ex}%
  {\normalfont\large\bfseries}}
\renewcommand{\@maketitle}{%
  \newpage\null\vskip -1.2em%
  {\footnotesize\bfseries
    \href{https://www.ufl.edu/}{University of Florida}, \href{https://math.ufl.edu/}{Department of Mathematics}\par}%
  \vskip 0.18em%
  \tagstructbegin{tag=Title}%
    {\LARGE\bfseries\href{https://gma.math.ufl.edu/exams/numerical-optimization-phd/}{Variational Analysis/Numerical Optimization PhD Exam}, May 2025\par}%
  \tagstructend%
  \vskip 0.35em%
  \tagmcbegin{artifact}%
    \hrule height 1pt%
  \tagmcend%
  \vskip 0.55em%
}
\makeatother
\title{Variational Analysis/Numerical Optimization PhD Exam, May 2025}
\author{}
\date{}
\begin{document}
\maketitle
\thispagestyle{empty}
\begin{examinstructions}
Do 8 problems, 4 from problems 1–5 and 4 from problems 6–10.
\end{examinstructions}
\begin{examproblems}{0}
\item
Recall that \(f:\mathbb{R}^n\to\mathbb{R}\) is convex if and only if 
\[
f[(1-\alpha)\mathbf{x}+\alpha\mathbf{y}]\leq (1-\alpha)f(\mathbf{x})+\alpha f(\mathbf{y})
\]
 for all \(\mathbf{x},\mathbf{y}\in\mathbb{R}^n\) and \(\alpha\in[0,1]\).
\begin{examsubparts}
\item[\textnormal{(a)}]
If~\(f\) is continuously differentiable, then show that~\(f\) is convex if and only if 
\[
f(\mathbf{y})\geq f(\mathbf{x})+\nabla f(\mathbf{x})(\mathbf{y}-\mathbf{x})
\]
 for all \(\mathbf{x},\mathbf{y}\in\mathbb{R}^n\).
\item[\textnormal{(b)}]
If~\(f\) is twice continuously differentiable, then show that~\(f\) is convex if and only if \(\nabla^2f\) is positive semidefinite on \(\mathbb{R}^n\).
\end{examsubparts}
\item
Suppose that \(\boldsymbol{\Phi}:\mathcal{K}\to\mathbb{R}^n\) is continuously differentiable on \(\mathcal{K}\subset\mathbb{R}^n\).
\begin{examsubparts}
\item[\textnormal{(a)}]
If~\(\mathcal{K}\) is a convex set, then show that 
\[
\|\boldsymbol{\Phi}(\mathbf{x})-\boldsymbol{\Phi}(\mathbf{y})\|\leq \mu\|\mathbf{x}-\mathbf{y}\|
\]
 for all~\(\mathbf{x}\) and \(\mathbf{y}\in\mathcal{K}\), where~\(\mu\) is the supremum of the singular values of \(\nabla\boldsymbol{\Phi}\) over~\(\mathcal{K}\).
\item[\textnormal{(b)}]
If \(\boldsymbol{\Phi}\) is a contraction on~\(\mathcal{K}\), where~\(\mathcal{K}\) is a closed set, and \(\boldsymbol{\Phi}(\mathbf{x})\in\mathcal{K}\) for each \(\mathbf{x}\in\mathcal{K}\), then show that \(\boldsymbol{\Phi}\) has a unique fixed point in~\(\mathcal{K}\).
\end{examsubparts}
\item
Let~\(\mathbf{U}\) and \(\mathbf{V}\in\mathbb{R}^{n\times m}\) and \(\mathbf{M}\in\mathbb{R}^{n\times n}\) with~\(\mathbf{M}\) invertible.
\begin{examsubparts}
\item[\textnormal{(a)}]
If \(\mathbf{I}+\mathbf{V}^{\mathsf T}\mathbf{M}^{-1}\mathbf{U}\) is invertible, then show that \(\mathbf{M}+\mathbf{U}\mathbf{V}^{\mathsf T}\) is invertible with 
\[
(\mathbf{M}+\mathbf{U}\mathbf{V}^{\mathsf T})^{-1}=\mathbf{M}^{-1}-\mathbf{M}^{-1}\mathbf{U}(\mathbf{I}+\mathbf{V}^{\mathsf T}\mathbf{M}^{-1}\mathbf{U})^{-1}\mathbf{V}^{\mathsf T}\mathbf{M}^{-1}.
\]
\item[\textnormal{(b)}]
If \(\mathbf{U}=\mathbf{V}=\mathbf{A}\in\mathbb{R}^{n\times m}\), where the columns of~\(\mathbf{A}\) are linearly independent, and \(\mathbf{M}=\mathbf{Q}\), a symmetric matrix, then show that 
\[
(\mathbf{Q}+p\mathbf{A}\mathbf{A}^{\mathsf T})(\mathbf{I}+p\mathbf{A}\mathbf{A}^{\mathsf T})^{-1}=\mathbf{Q}_0+\mathcal{O}(1/p),
\]
 where 
\[
\mathbf{Q}_0=\mathbf{Q}+(\mathbf{I}-\mathbf{Q})[\mathbf{A}(\mathbf{A}^{\mathsf T}\mathbf{A})^{-1}\mathbf{A}^{\mathsf T}].
\]
\end{examsubparts}
\item
Consider the following quadratic program: 
\[
\min\ q(\mathbf{x}):=\sum_{i=1}^n 0.5x_i^2+c_i x_i
\quad\text{subject to}\quad \mathbf{a}^{\mathsf T}\mathbf{x}=b,\quad \mathbf{x}\geq 0. \tag{QP}
\]
 where~\(\mathbf{c}\) and \(\mathbf{a}\in\mathbb{R}^n\), \(\mathbf{a}\geq 0\), and \(b>0\) is a scalar.
\begin{examsubparts}
\item[\textnormal{(a)}]
Develop an algorithm for solving (QP) by maximizing the dual function 
\[
L(\lambda)=\inf\{q(\mathbf{x})+\lambda(\mathbf{a}^{\mathsf T}\mathbf{x}-b):\mathbf{x}\geq 0\}
\]
 over the dual multiplier~\(\lambda\). Here, the constraint~\(\mathcal{K}\) in the dual function is taken as \(\mathcal{K}=\{\mathbf{x}\in\mathbb{R}^n:\mathbf{x}\geq 0\}\).
\item[\textnormal{(b)}]
Given the optimal~\(\lambda\) for the dual problem, what is the solution of the primal problem?
\end{examsubparts}
\item
Consider the primal problem 
\[
\min\ f(\mathbf{x})\quad\text{subject to}\quad \mathbf{A}\mathbf{x}-\mathbf{b}=0,\quad \mathbf{g}(\mathbf{x})\leq 0, \tag{P}
\]
 where \(f:\mathbb{R}^n\to\mathbb{R}\) and the components of \(\mathbf{g}:\mathbb{R}^n\to\mathbb{R}^l\) are convex and continuously differentiable, \(\mathbf{A}\in\mathbb{R}^{m\times n}\), and \(\mathbf{b}\in\mathbb{R}^m\). The dual problem is 
\[
\max\ L(\boldsymbol{\lambda},\boldsymbol{\mu})\quad\text{subject to}\quad \boldsymbol{\lambda}\in\mathbb{R}^m,\quad \boldsymbol{\mu}\in\mathbb{R}^l,\quad \boldsymbol{\mu}\geq 0, \tag{D}
\]
 where 
\[
\begin{aligned}
L(\boldsymbol{\lambda},\boldsymbol{\mu})&=\inf\{\mathcal{L}(\mathbf{x},\boldsymbol{\lambda},\boldsymbol{\mu}):\mathbf{x}\in\mathbb{R}^n\},\\
\mathcal{L}(\mathbf{x},\boldsymbol{\lambda},\boldsymbol{\mu})&=f(\mathbf{x})+\boldsymbol{\lambda}^{\mathsf T}(\mathbf{A}\mathbf{x}-\mathbf{b})+\boldsymbol{\mu}^{\mathsf T}\mathbf{g}(\mathbf{x}).
\end{aligned}
\]
\begin{examsubparts}
\item[\textnormal{(a)}]
Show that if (P) has a local minimizer \(\mathbf{x}^*\), then \(\mathbf{x}^*\) is a global minimizer for (P). Moreover, if the first-order optimality conditions hold at \(\mathbf{x}^*\), then there is a solution to (D) with no duality gap.
\item[\textnormal{(b)}]
Returning to (QP) of problem 4, show that there is no duality gap. You may wish to consider both the set 
\[
S=\{(r,s)\in\mathbb{R}^2:r\geq q(\mathbf{x})\text{ and }s=\mathbf{a}^{\mathsf T}\mathbf{x}-b\text{ for some }\mathbf{x}\geq 0\},
\]
 and the point \((q^*,0)\), where \(q^*\) is the optimal cost for (QP). Explain why this point does not lie in the interior of~\(S\). Use the separating hyperplane theorem to separate~\(S\) and the point, and then analyze the resulting inequality.
\end{examsubparts}
\item
Consider the initial-value problem 
\[
\dot{\mathbf{x}}(t)=\mathbf{f}(\mathbf{x}(t),\mathbf{u}(t)),\qquad \mathbf{x}(0)=\mathbf{a},
\]
 where \(\mathbf{a}\in\mathbb{R}^n\), \(\mathbf{f}:\mathcal{B}_\delta(\mathbf{a})\times\mathcal{U}\to\mathbb{R}^n\), with \(\mathcal{U}\subset\mathbb{R}^m\) a compact set and \(\delta>0\). It is assumed that~\(\mathbf{f}\) is continuous on \(\mathcal{B}_\delta(\mathbf{a})\times\mathcal{U}\) and uniformly Lipschitz continuous in its first argument with Lipschitz constant~\(L\) satisfying 
\[
|\mathbf{f}(\mathbf{x}_1,\mathbf{u})-\mathbf{f}(\mathbf{x}_2,\mathbf{u})|\leq L|\mathbf{x}_1-\mathbf{x}_2|
\]
 for all \(\mathbf{x}_1,\mathbf{x}_2\in\mathcal{B}_\delta(\mathbf{a})\) and \(\mathbf{u}\in\mathcal{U}\). If \(\epsilon=\delta/M\), where 
\[
M=\max\{|\mathbf{f}(\mathbf{x},\mathbf{u})|:\mathbf{x}\in\mathcal{B}_\delta(\mathbf{a}),\ \mathbf{u}\in\mathcal{U}\},
\]
 then show that the initial-value problem has a unique, absolutely continuous solution on the interval \([0,\epsilon]\) with \(\mathbf{x}(t)\in\mathcal{B}_\delta(\mathbf{a})\) for all \(t\in[0,\epsilon]\).
\item
Consider the equation 
\[
(P(x)u'(x))'=Q(x)u(x),\qquad x\in[0,1]. \tag{E}
\]
 where~\(Q\) is continuous and~\(P\) is continuously differentiable and strictly positive on \([0,1]\). Show that if the solution to the initial-value problem for (E) with \(u(0)=0\) and \(u'(0)=1\) is strictly positive on the half-open interval \((0,1]\), then there exists a strictly positive solution to the second-order differential equation (E) (without boundary conditions at \(x=0\) or \(x=1\)).
\item
Consider the variational problem 
\[
\min\ \int_0^1\left(\frac{1}{2}u'(x)^2+e^{u(x)}\right)\,dx
\quad\text{subject to}\quad u\in\mathcal{H}_0^1.
\]
\begin{examsubparts}
\item[\textnormal{(a)}]
What is the first-order necessary optimality condition (Euler equation) for this variational problem?
\item[\textnormal{(b)}]
Consider a uniform mesh \(x_k=kh\), where \(h=1/N\); let \(u_k\) denote an approximation to \(u(x_k)\). Of course, \(u_0=u_N=0\). Give a finite difference approximation in terms \(u_1,\ldots,u_{N-1}\) to the solution of the Euler equation.
\item[\textnormal{(c)}]
Let \(\mathbf{F}(\mathbf{u})\) denote the finite difference system where \(\mathbf{u}=(u_1,u_2,\ldots,u_{N-1})^{\mathsf T}\in\mathbb{R}^{N-1}\), and let \(\mathbf{u}^*\) denote the vector formed by evaluating the solution of the Euler equation at the mesh points \(x_1,x_2,\ldots,x_{N-1}\). Obtain a bound for the components of \(\mathbf{F}(\mathbf{u}^*)\) in terms of the mesh spacing~\(h\).
\end{examsubparts}
\item
Consider the initial-value problem 
\[
\dot{\mathbf{x}}(t)=\mathbf{A}(t)\mathbf{x}(t)+\mathbf{u}(t),\qquad \mathbf{x}(0)=0,
\]
 where \(\mathbf{x}:[0,1]\to\mathbb{R}^n\) and \(\mathbf{A}\in\mathcal{L}^{\infty}([0,1];\mathbb{R}^{n\times n})\).
\begin{examsubparts}
\item[\textnormal{(a)}]
Assuming \(\mathbf{u}\in\mathcal{L}^2\), obtain a bound for the sup-norm of~\(\mathbf{x}\) in terms of the \(\mathcal{L}^2\) norm of~\(\mathbf{u}\).
\item[\textnormal{(b)}]
Let \(\Omega:\mathcal{H}_0^1\to\mathbb{R}\) be defined by 
\[
\Omega(h)=\frac{1}{2}\int_0^1 r^2(x)\,dx\qquad\text{where }r(x)=h'(x)+w(x)h(x).
\]
 Show that there exists \(\beta>0\) such that 
\[
\Omega(h)\geq\beta(\|h\|^2+\|h'\|^2)\qquad\text{for all }h\in\mathcal{H}_0^1,
\]
 where \(\|\cdot\|\) is the \(\mathcal{L}^2\) norm.
\item[\textnormal{(c)}]
Consider the quadratic programming problem 
\[
\min\ \frac{1}{2}\mathbf{x}^{\mathsf T}\mathbf{R}\mathbf{x}+\mathbf{b}^{\mathsf T}\mathbf{x}
\quad\text{subject to}\quad \mathbf{x}\in\mathcal{K},
\]
 where \(\mathcal{K}\subset\mathbb{R}^n\) is a closed convex set and \(\mathbf{R}\in\mathbb{R}^{n\times n}\) is positive definite with smallest eigenvalue \(\alpha>0\). Show that the solutions \(\mathbf{u}_i\), \(i=1,2\), associated with \(\mathbf{b}=\mathbf{b}_i\), \(i=1,2\) respectively, satisfy 
\[
\|\mathbf{u}_1-\mathbf{u}_2\|\leq \|\mathbf{b}_1-\mathbf{b}_2\|/\alpha.
\]
\end{examsubparts}
\item
Consider the following control problem: 
\[
\min\ \int_0^1\left(\frac{1}{2}u^2(x)+y(x)\right)\,dx
\quad\text{subject to}\quad y'(x)=u(x),\quad y(0)=0,\quad u(x)\geq\ell(x),
\]
 where \(\ell(x)\) is a given lower bound for the control \(u(x)\) at each \(x\in[0,1]\).
\begin{examsubparts}
\item[\textnormal{(a)}]
What is the first-order optimality condition (Pontryagin minimum principle) for this problem?
\item[\textnormal{(b)}]
What is the solution of the control problem?
\end{examsubparts}
\end{examproblems}
\end{document}
